\documentclass[11pt]{article}
\usepackage[margin=1in]{geometry}
\usepackage[T1]{fontenc}
\usepackage[utf8]{inputenc}
\usepackage{amsmath,amssymb,amsthm}
\usepackage{graphicx}
\usepackage{xcolor}
\usepackage{microtype}
\usepackage[colorlinks=true,linkcolor=blue!55!black,citecolor=blue!55!black,urlcolor=blue!55!black]{hyperref}

\newtheorem{theorem}{Theorem}
\newtheorem{lemma}{Lemma}
\newcommand{\Bop}{\mathcal B}
\newcommand{\Kop}{\mathcal K}

\title{No proper operator ideal has the Berger--Coburn phenomenon on Hardy or Bergman space}
\author{Candidate full-solution packet for arXiv:2304.09716, Open Problem 1}
\date{11 August 2026}

\begin{document}
\maketitle

\begin{center}
\fbox{\parbox{0.92\textwidth}{\textbf{Status: candidate full solution, likely valid, subject to human review.}
Under the standard meaning of a proper two-sided operator ideal, the answer to
Open Problem~1 is no.  Every proper ideal between infinite-dimensional Hilbert
spaces consists of compact operators.  The source paper's analytic-symbol
examples therefore rule out the Berger--Coburn phenomenon for every proper
ideal at once.}}
\end{center}

\section{The open problem}

Let $X$ be the Hardy space $H^2(\mathbb T)$ or the Bergman space
$A^2(\mathbb B_n)$, let $Y$ be the corresponding $L^2$ space, and let
\[
 H_f=(I-P)M_f:X\longrightarrow Y
\]
be the Hankel operator with bounded symbol $f$.  The source says that an
operator ideal $\mathcal M\subset\Bop(X,Y)$ has the Berger--Coburn phenomenon
(BCP) when
\[
 H_f\in\mathcal M\quad\Longleftrightarrow\quad H_{\bar f}\in\mathcal M
 \qquad(f\in L^\infty).
\]
After exhibiting failure for compact operators, Hu and Virtanen ask
\cite[Open Problem~1]{source}:
\begin{quote}
Are there nontrivial operator ideals in
$\Bop(H^2(\mathbb T),L^2(\mathbb T))$ or in
$\Bop(A^2(\mathbb B_n),L^2(\mathbb B_n))$ for which the
Berger--Coburn phenomenon holds?
\end{quote}
Here ``nontrivial'' is read in the standard sense ``proper,'' excluding the
full class $\Bop(X,Y)$, for which membership is vacuous.  The answer is no.

Figure~\ref{fig:source} reproduces printed page~4, including both the decisive
analytic-symbol examples and Open Problem~1.

\begin{figure}[ht]
\centering
\includegraphics[width=0.94\textwidth]{figures/open_problem_crop.png}
\caption{The source counterexamples and Open Problem~1.}
\label{fig:source}
\end{figure}

\clearpage
\section{Idea of the proof}

There are only two ingredients.  First, a proper two-sided operator ideal
between infinite-dimensional Hilbert spaces cannot contain a noncompact
operator.  Indeed, every noncompact $T$ is bounded below on some
infinite-dimensional closed subspace.  After identifying the domain with that
subspace, $T$ has a bounded left inverse, so the identity factors through $T$.
The ideal property then puts every bounded operator in the ideal.

Second, the source paper already gives a bounded analytic symbol $b$ in each
Hardy/Bergman setting such that $H_b=0$ while $H_{\bar b}$ is noncompact.  The
zero operator belongs to every ideal, whereas the lemma excludes
$H_{\bar b}$ from every proper ideal.  Thus the BCP equivalence fails for all
proper ideals simultaneously.

\section{The Hilbert-space ideal lemma}

For clarity, by an operator ideal in $\Bop(X,Y)$ we mean a linear subspace
$\mathcal M$ satisfying
\[
 A T B\in\mathcal M
 \quad(T\in\mathcal M,\ A\in\Bop(Y),\ B\in\Bop(X)).
\]
This is exactly the rectangular component of a standard two-sided operator
ideal.

\begin{lemma}\label{lem:ideal}
Let $X,Y$ be separable infinite-dimensional Hilbert spaces.  If
$\mathcal M\subsetneq\Bop(X,Y)$ is an operator ideal, then
\[
 \mathcal M\subseteq\Kop(X,Y).
\]
No norm-closure hypothesis is required.
\end{lemma}

\begin{proof}
Suppose $T\in\mathcal M$ is not compact.  Write $|T|=(T^*T)^{1/2}$.  There
is an $\varepsilon>0$ for which the spectral subspace
\[
 E=\mathbf 1_{[\varepsilon,\infty)}(|T|)X
\]
is infinite-dimensional.  Otherwise all these spectral subspaces would be
finite-dimensional, and the finite-rank operators
$|T|\mathbf 1_{[1/m,\infty)}(|T|)$ would converge in norm to $|T|$, making
$T$ compact.

Choose an isometry $V:X\to X$ with range $E$.  Then
\[
 \|TVx\|=\||T|Vx\|\geq\varepsilon\|x\|,
\]
so $TV$ has closed range and a bounded left inverse.  More explicitly, define
$S:Y\to X$ by $S(TVx)=x$ on $\operatorname{ran}(TV)$ and by zero on its
orthogonal complement.  Then $STV=I_X$.

For arbitrary $R\in\Bop(X,Y)$,
\[
 R=(RS)TV.
\]
Since $T\in\mathcal M$ and $RS\in\Bop(Y)$, $V\in\Bop(X)$, the ideal property
gives $R\in\mathcal M$.  Hence $\mathcal M=\Bop(X,Y)$, contradicting
properness.  Every member of a proper ideal is therefore compact.
\end{proof}

\section{Negative answer to Open Problem 1}

\begin{theorem}\label{thm:main}
Let
\[
 (X,Y)=\bigl(H^2(\mathbb T),L^2(\mathbb T)\bigr)
 \quad\text{or}\quad
 (X,Y)=\bigl(A^2(\mathbb B_n),L^2(\mathbb B_n)\bigr).
\]
No proper operator ideal $\mathcal M\subsetneq\Bop(X,Y)$ has the
Berger--Coburn phenomenon for all bounded symbols.  Thus the answer to
Open Problem~1 is negative in both settings.
\end{theorem}

\begin{proof}
All four Hilbert spaces in the statement are separable and
infinite-dimensional.  Lemma~\ref{lem:ideal} gives
$\mathcal M\subseteq\Kop(X,Y)$ for every proper ideal $\mathcal M$.

For the Hardy space, the source chooses
$b\in H^\infty\setminus\mathrm{VMO}$; then $H_b=0$ but
$H_{\bar b}$ is not compact.  For the Bergman space on the disk, the source
uses an explicit infinite Blaschke product with the same two properties, and
it extends this example by $b(z)=b(z_1)$ to $\mathbb B_n$
\cite[Example~2]{source}.  In every case,
\[
 H_b=0\in\mathcal M,
 \qquad
 H_{\bar b}\notin\Kop(X,Y)\supseteq\mathcal M.
\]
Therefore the defining equivalence for BCP fails.  Since $\mathcal M$ was
arbitrary, no proper operator ideal works.
\end{proof}

\section{Scope, novelty check, and review recommendation}

The theorem answers exactly Open Problem~1 under the standard two-sided ideal
meaning used for compact and Schatten ideals.  If ``ideal'' were intended to
mean a symbol-dependent class lacking stability under arbitrary left and right
bounded multiplication, Lemma~\ref{lem:ideal} would not apply; such a class is
not an operator ideal in the usual sense.  The full class $\Bop(X,Y)$ has BCP
trivially and is precisely the excluded trivial case.

On 11 August 2026, the run registry and all cheap indexes had no prior answer
for arXiv:2304.09716 or this ideal-factorization route.  Exact-title, exact-
problem, author, Hardy/Bergman operator-ideal, and Berger--Coburn searches of
the local arXiv corpus and the web found the source and later work on doubling
Fock spaces and symmetrically normed ideals, but no answer to Open Problem~1.
The later doubling-Fock paper arXiv:2312.06656 explicitly makes partial
progress on a different open problem from the source.  This is a bounded
novelty check, not an exhaustive priority claim.

Human review should focus on the intended meaning of ``operator ideal'' in
Open Problem~1.  Once the standard two-sided meaning is confirmed, the
factorization in Lemma~\ref{lem:ideal} and the source's noncompactness examples
make the conclusion immediate.

\begin{thebibliography}{9}
\bibitem{source}
Z. Hu and J. A. Virtanen,
\emph{On the Berger--Coburn phenomenon},
Operator Theory: Advances and Applications 297 (2024), 223--234;
arXiv:2304.09716.
\end{thebibliography}

\end{document}
